\section*{Hoofdstuk \ref{ch:b2:solid-liquid-diagrams} --- Binaire diagrammen vast-vloeibaar}

\begin{solution}{exo:b2:solid-liquid-diagrams:1}
Bij \qty{1300}{\degreeCelsius} ligt de \omterm{def:b2:solid-liquid-diagrams:liquidus}{liquidus} bij 0.541 en de \omterm{def:b2:solid-liquid-diagrams:liquidus}{solidus} bij 0.609.
0.40: alleen vloeistof. 0.58: vloeistof (0.541) en \omterm{def:b2:solid-liquid-diagrams:solid-solution}{vaste oplossing} (0.609). 0.70:
alleen \omterm{def:b2:solid-liquid-diagrams:solid-solution}{vaste oplossing}.
\end{solution}

\begin{solution}{exo:b2:solid-liquid-diagrams:2}
Fractie vaste stof $(0.58 - 0.541)/(0.609 - 0.541) = 0.57$: \qty{1.15}{mol} vaste stof,
\qty{0.85}{mol} vloeistof. Nikkel: $1.15 \times 0.609 = \qty{0.70}{mol}$ in de
vaste stof, $0.85 \times 0.541 = \qty{0.46}{mol}$ in de vloeistof; totaal \qty{1.16}{mol}
$= 2.00 \times 0.58$.
\end{solution}

\begin{solution}{exo:b2:solid-liquid-diagrams:3}
Bij vaste druk is $v = n + 1 - \varphi$ (geen reactie). (a) $1 + 1 - 2 = 0$.
(b) $2 + 1 - 2 = 1$. (c) Vloeistof en twee \omterm{def:b2:solid-liquid-diagrams:solid-solution}{vaste oplossingen}: $2 + 1 - 3 = 0$.
\end{solution}

\begin{solution}{exo:b2:solid-liquid-diagrams:4}
Een gelijkmatige daling tot \qty{232}{\degreeCelsius}, een horizontaal plateau terwijl het tin
stolt, en daarna opnieuw een gelijkmatige daling, trager dicht bij kamertemperatuur. Op het
plateau bestaan vloeistof en vaste stof naast elkaar bij vaste druk, \omterm{def:b2:equilibrium-shifts:variance}{variantie} 0: de
afgevoerde warmte wordt geleverd door de kristallisatie, en de temperatuur kan niet veranderen
tot de laatste vloeistof gestold is.
\end{solution}

\begin{solution}{exo:b2:solid-liquid-diagrams:5}
Benzeen: $\ln x_B = -(9870/8.314)(1/269.6 - 1/278.64) = -0.1429$, $x_B =
0.867$. Naftaleen: $\ln x_N = -(19\,100/8.314)(1/269.6 - 1/353.2) = -2.017$,
$x_N = 0.133$. De som is 1.000: \qty{269.6}{K} ($\qty{-3.5}{\degreeCelsius}$) is
de eutectische temperatuur, en de eutectische vloeistof bevat 0.133 naftaleen.
\end{solution}

\begin{solution}{exo:b2:solid-liquid-diagrams:6}
De vloeistof koelt af tot de \omterm{def:b2:solid-liquid-diagrams:liquidus}{liquidus} aan de loodkant, waar kristallen (Pb) verschijnen; de
vloeistof wordt rijker aan tin, tot aan het eutecticum bij \qty{182.2}{\degreeCelsius}
en 62.1~\%. Daar stolt de resterende vloeistof tot het eutectische mengsel.
Massafractie eutecticum: $(30 - 18.91)/(62.13 - 18.91) = 0.26$. Bij verdere
afkoeling daalt de oplosbaarheid van tin in (Pb) (gestreepte lijn): een beetje (Sn)
slaat neer in de kristallen (Pb).
\end{solution}

\begin{solution}{exo:b2:solid-liquid-diagrams:7}
23.3~\% NaCl in massa: $23.3/76.7 = \qty{0.304}{kg}$ zout per kilogram
water. Bij \qty{-25}{\degreeCelsius}, onder het eutecticum, kan bij geen enkele verhouding
een vloeistof bestaan: zout en ijs blijven als vaste stoffen naast elkaar liggen.
\end{solution}

\begin{solution}{exo:b2:solid-liquid-diagrams:8}
$x_{\mathrm B} = 0.5$ ligt tussen E$_1$ en de verbinding: $\mathrm{AB_2}$
kristalliseert eerst, de vloeistof schuift naar E$_1$ en stolt daar tot het
eutecticum van A en $\mathrm{AB_2}$; op het einde A + $\mathrm{AB_2}$.
$x_{\mathrm B} = 0.9$ ligt tussen E$_2$ en B: B kristalliseert eerst, de vloeistof
bereikt E$_2$; op het einde $\mathrm{AB_2}$ + B.
\end{solution}

\begin{solution}{exo:b2:solid-liquid-diagrams:9}
Aan de achterrand van de zone stolt de vloeistof tot een vaste stof met slechts 0.78
maal haar kopermolfractie: het afgestoten koper blijft in de vloeistof, die
verder schuift en rijker wordt, zodat het koper naar het uiteinde van de staaf wordt geveegd.
Met een verhouding dicht bij 1 verwijdert elke doorgang slechts een fractie van de onzuiverheid:
er zijn veel doorgangen nodig, elk vertrekkend van de staaf die de vorige achterliet.
\end{solution}

\begin{solution}{exo:b2:solid-liquid-diagrams:10}
Afleiding zoals in het hoofdstuk. Voor een verdunde vloeistof is $\ln x_A \approx -x_B$, $1/T -
1/T^* \approx -\Delta T/T^{*2}$, dus $x_B = \Delta_{\text{fus}}H_A\Delta T/(RT^{*2})$;
met $x_B \approx n_B/n_A = bM_A$ is $\Delta T = (RT^{*2}M_A/\Delta_{\text{fus}}H_A)\,b$.
Benzeen: $K_f = 8.314 \times 278.64^2 \times 0.07811/9870 = \qty{5.11}{K.kg/mol}$.
\end{solution}

\begin{solution}{exo:b2:solid-liquid-diagrams:11}
Per \qty{100}{g}: $19.0/24.31 = \qty{0.782}{mol}$ Mg, $81.0/207.2 =
\qty{0.391}{mol}$ Pb; verhouding $2.00$: $\ce{Mg2Pb}$.
\end{solution}

\begin{solution}{exo:b2:solid-liquid-diagrams:12}
Het plateau is evenredig met de massa eutectische vloeistof: de onbekende bevat
er $120/300 = 0.40$ van. Aan de loodkant is $w = 18.91 + 0.40 \times (62.13 -
18.91) = 36.2~\%$ tin; aan de tinkant $w = 96.59 - 0.40 \times (96.59 - 62.13) =
82.8~\%$. De eerste knik van de \omterm{def:b2:solid-liquid-diagrams:cooling-curve}{afkoelingskromme} (hoger voor de legering met 36~\%,
waarvan de \omterm{def:b2:solid-liquid-diagrams:liquidus}{liquidus} aan de steile loodkant ligt) of een microfoto (primaire kristallen (Pb) of
primaire kristallen (Sn)) onderscheidt ze.
\end{solution}

\begin{solution}{pb:b2:solid-liquid-diagrams:1}
\textbf{1.} Lood \qty{327.5}{\degreeCelsius}, tin \qty{231.9}{\degreeCelsius}.
\textbf{2.} Per \qty{100}{g}: $62.13/118.71 = \qty{0.5234}{mol}$ tin en
$37.87/207.2 = \qty{0.1828}{mol}$ lood; $x_{\mathrm{Sn}} = 0.5234/0.7062 =
0.741$.
\textbf{3.} Punten (0, 327.5), (100, 231.9), (62.13, 182.2), (18.91, 182.2),
(96.59, 182.2).
\textbf{4.} Zoals in de figuur: vloeistof; L + (Pb); L + (Sn); (Pb); (Sn); (Pb) + (Sn).
\textbf{5.} L (62.1~\% Sn) $\longrightarrow$ (Pb) (18.9~\%) + (Sn) (96.6~\%).
\textbf{6.} $v = 2 + 1 - 3 = 0$: de temperatuur en de drie samenstellingen liggen
vast.
\textbf{7.} De loodrijke \omterm{def:b2:solid-liquid-diagrams:solid-solution}{vaste oplossing} (Pb): lood lost tin op, dus de vaste stof in
evenwicht met de vloeistof bevat al tin, in de verhouding die op de
\omterm{def:b2:solid-liquid-diagrams:liquidus}{solidus} wordt afgelezen.
\textbf{8.} Ze volgt de \omterm{def:b2:solid-liquid-diagrams:liquidus}{liquidus} naar het eutecticum en wordt rijker aan tin,
tot 62.13~\%.
\textbf{9.} Vloeistof met 62.13~\% en (Pb) met 18.91~\% tin.
\textbf{10.} (Pb): $(62.13 - 40)/(62.13 - 18.91) = 0.512$; vloeistof 0.488.
\textbf{11.} (Pb) met 18.91~\% en (Sn) met 96.59~\%; (Sn): $(40 - 18.91)/(96.59 -
18.91) = 0.271$; (Pb): 0.729.
\textbf{12.} Eutectisch bestanddeel 48.8~\% (de vloeistof die er net boven aanwezig is);
primair (Pb) 51.2~\%. De (Pb) van vraag 11 is de som van de primaire
kristallen en van de lamellen (Pb) van het eutecticum.
\textbf{13.} Een knik op de \omterm{def:b2:solid-liquid-diagrams:liquidus}{liquidus}, een tragere daling, dan een plateau bij
\qty{182.2}{\degreeCelsius} dat 0.488 maal zo lang duurt als dat van de eutectische legering.
\textbf{14.} $\ln x_{\mathrm{Sn}} = -(7030/8.314)(1/T - 1/505.1)$.
\textbf{15.} $\ln x_{\mathrm{Bi}} = -(11\,300/8.314)(1/T - 1/544.6)$.
\textbf{16.} $x_{\mathrm{Sn}}(T) + x_{\mathrm{Bi}}(T) = 1$.
\textbf{17.} Bij \qty{110}{\degreeCelsius}: $0.587 + 0.349 = 0.936$; bij
\qty{120}{\degreeCelsius}: $0.621 + 0.382 = 1.003$; bij
\qty{130}{\degreeCelsius}: $0.655 + 0.417 = 1.072$. De som bereikt 1 rond
\qty{119.5}{\degreeCelsius}: \qty{120}{\degreeCelsius} op een graad nauwkeurig.
\textbf{18.} $x_{\mathrm{Bi}} = 0.38$; in massa $0.38 \times 208.98/(0.38 \times
208.98 + 0.62 \times 118.71) = 0.52$.
\textbf{19.} Het model geeft \qty{120}{\degreeCelsius} en 52~\% bismut, het
geëvalueerde diagram \qty{138.8}{\degreeCelsius} en 57~\%: ongeveer
\qty{19}{\degreeCelsius} te laag.
\textbf{20.} Het tin dat kristalliseert, is niet zuiver maar een oplossing die
bismut bevat: zijn \omterm{def:b2:chemical-potential:chemical-potential}{chemische potentiaal} is lager dan die van zuiver tin, de vaste stof is
stabieler, en ze kristalliseert bij een gegeven samenstelling van de vloeistof bij een hogere
temperatuur. De tintak van de \omterm{def:b2:solid-liquid-diagrams:liquidus}{liquidus} wordt verhoogd en snijdt de bismuttak
hoger (een niet-ideale vloeistof voegt er haar eigen correctie aan toe).
\textbf{21.} Ze smelt en stolt bij één temperatuur, de laagste van het systeem:
geen pasteus gebied waarin een verbinding die beweegt, zou scheuren.
\textbf{22.} Lood is giftig, en elektronisch afval komt in het milieu terecht.
\textbf{23.} Warmtegevoelige componenten kunnen bij een lagere temperatuur gesoldeerd worden;
maar de verbinding wordt in gebruik al bij een lagere temperatuur zacht, en bismut maakt haar
bros.
\textbf{24.} \qty{570}{g} bismut en \qty{430}{g} tin.
\textbf{25.} Het model van de ideale vloeistof voorspelt een eutecticum bij $\boldsymbol{T_E
\approx \qty{120}{\degreeCelsius}}$, tegenover \qty{138.8}{\degreeCelsius}
geëvalueerd.
\end{solution}
